\chapter{Implementation and Validation}

This chapter presents the implementation of the proposed PCAP preprocessing system. While the previous chapter described the overall architecture and information flow, this chapter focuses on the practical development of the different software components.

The implementation follows a modular approach in which the agent orchestration, specialized tools, and underlying processing logic are separated according to their responsibilities. The system is organized around three main preprocessing operations: automotive event extraction, noise filtering and context compression, and context building and report generation.

The implementation also includes the agent workflow responsible for coordinating these operations and maintaining the information required during execution.

\section{Development Methodology}

The development of the system followed an incremental implementation approach. The preprocessing workflow was divided into smaller functional components, allowing each processing stage to be developed and validated independently before being integrated into the complete agent workflow.

The development process began with the implementation of the core processing capabilities required to extract relevant information from communication data. The filtering and context compression functionality was then introduced to reduce repetitive or low-value information while preserving meaningful communication context.

Finally, the outputs generated by the preprocessing stages were aggregated into a unified report intended to provide a structured input for downstream AI analysis.

This incremental approach simplified the development and validation process. Each component could be tested individually, making it easier to identify implementation issues before executing the complete workflow.

\section{Project Implementation Structure}

\subsection{Separation of Responsibilities}

The implementation was organized according to a separation between three main responsibilities:

\begin{itemize}
    \item agent and workflow orchestration;
    \item specialized preprocessing tools;
    \item core data processing logic.
\end{itemize}

This organization was adopted to maintain a clear distinction between workflow coordination and the actual processing of communication data.

The agent layer is responsible for coordinating the execution of the workflow. The tool layer provides the operations that can be invoked by the agent. The core processing layer contains the implementation responsible for performing the underlying extraction, filtering, compression, and report generation operations.

This separation improves maintainability because modifications to the processing logic can be performed without unnecessarily modifying the agent workflow.

\subsection{Core Processing Layer}

The core processing layer contains the main implementation logic of the preprocessing system.

This layer is responsible for performing the actual data transformations required by the different stages of the pipeline. These operations include processing extracted communication information, identifying repetitive or low-value data, preserving relevant context, and preparing the information required for the final report.

Keeping this logic within a dedicated processing layer avoids placing complex processing operations directly inside the agent workflow or tool definitions.

The core processing layer therefore represents the main functional implementation of the preprocessing system.

\subsection{Tool Layer}

The tool layer provides the interface between the agent workflow and the underlying processing operations.

Each tool corresponds to a specific preprocessing capability. Rather than containing unnecessary orchestration logic, the tools provide a clear entry point for invoking the corresponding processing functionality.

The implemented preprocessing tools are:

\begin{enumerate}
    \item Automotive Event Extraction;
    \item Noise Filtering and Context Compression;
    \item Context Builder and Report Generation.
\end{enumerate}

The output produced by one preprocessing stage can be used as the input of the next stage, allowing the communication information to be progressively transformed throughout the workflow.

\subsection{Agent and Workflow Layer}

The agent and workflow layer is responsible for coordinating the execution of the preprocessing system.

The workflow defines the processing sequence and manages the interaction between the agent and the available tools. It also maintains the execution state required to preserve information throughout the different stages.

The workflow architecture is designed to avoid embedding detailed packet-processing operations within the agent itself. Instead, specialized tools perform the required preprocessing operations while the workflow coordinates their execution.

\section{Preprocessing Tool Implementation}

\subsection{Automotive Event Extraction}

The first preprocessing tool is responsible for extracting meaningful communication events from the input data.

Its purpose is to provide an initial structured representation of the communication activity that can be processed by subsequent stages. Instead of requiring later stages to operate directly on the complete raw communication trace, the extraction stage identifies information relevant to the preprocessing workflow.

The extracted output acts as the input of the noise filtering and context compression stage.

This separation is important because the second tool does not need to perform the initial event extraction again. Its responsibility is to process and refine the communication information already produced by the first stage.

\subsection{Noise Filtering and Context Compression}

The second tool receives the output generated by the automotive event extraction stage.

The objective of this component is to reduce unnecessary communication information while preserving the information required to maintain an accurate representation of relevant network behavior.

The implementation focuses on identifying information that provides limited additional analytical value when repeated. Depending on the available communication data, this can include repetitive events, duplicate information, keep-alive communication, broadcast traffic, or management-related traffic.

However, the filtering operation is not implemented as a simple process of removing all repeated communication. Repetition can sometimes provide useful information regarding communication behavior, timing, or changes in the state of a network interaction.

For this reason, the filtering process is designed to reduce unnecessary repetition while preserving the communication context required by the next stage of the preprocessing workflow.

The resulting output represents a reduced and more manageable communication context.

\subsection{Context Builder and Report Generation}

The third preprocessing tool is responsible for building the final communication context.

This component receives the information generated by the previous preprocessing stages and aggregates it into a standardized output.

The report generation process organizes the available information into a consistent structure that can be consumed by a downstream analysis system. The generated context can include relevant communication metadata, extracted events, information related to communication flows, detected protocol information, timing-related information, anomalies when available, and summary statistics.

The main purpose of this component is to establish a clear boundary between the preprocessing workflow and the future AI analysis layer.

As a result, the downstream analysis system does not need to directly process the complete raw PCAP trace. Instead, it receives the structured communication context generated by the preprocessing pipeline.

\section{AI Agent Workflow Implementation}

\subsection{Shared State}

The workflow uses a shared state to maintain information during execution.

The state acts as a common structure through which the different processing nodes access and update the information required by the workflow.

The stored information can include the processing request, agent messages, input information, intermediate tool outputs, and the final generated communication context.

Using a shared state provides a structured mechanism for exchanging information between processing stages. Each node receives the current state, performs its corresponding operation, and returns the updated information.

This approach avoids the use of unrelated global variables and allows the workflow to maintain a clear information flow.

\subsection{Processing Nodes}

The agent workflow is composed of nodes representing the main execution stages.

Each node receives the current workflow state and performs a specific responsibility. Depending on the node, this can include interacting with the language model, invoking a preprocessing tool, or updating the information maintained by the workflow.

The use of dedicated nodes provides a modular representation of the execution process. Individual components can be modified or extended without requiring changes to the complete system.

The specialized preprocessing tools are integrated into the workflow through their corresponding processing nodes.

\subsection{Workflow Execution}

The execution of the workflow begins when the input communication data and processing request are provided to the system.

The agent workflow coordinates the available processing operations and maintains the information generated during execution.

The logical preprocessing sequence follows the progression:

\begin{center}
\textbf{Input Communication Data}
\[
\Downarrow
\]
\textbf{Automotive Event Extraction}
\[
\Downarrow
\]
\textbf{Noise Filtering and Context Compression}
\[
\Downarrow
\]
\textbf{Context Builder and Report Generation}
\[
\Downarrow
\]
\textbf{Structured Communication Context}
\end{center}

The final result represents the output of the complete preprocessing workflow and is intended to be used by a downstream AI analysis component.

\section{Testing and Validation}

\subsection{Individual Tool Testing}

The preprocessing tools were validated individually to verify that each component could execute its intended operation.

Testing the components independently allowed the outputs of each preprocessing stage to be inspected before integration into the complete workflow.

The first tool was validated by verifying that relevant communication events could be extracted from the input data.

The filtering and context compression tool was then validated using the extracted output. The objective of this validation was to verify that unnecessary or repetitive information could be reduced while preserving meaningful communication information.

Finally, the context builder and report generator was tested to verify that the available preprocessing outputs could be aggregated into a consistent final result.

\subsection{End-to-End Workflow Testing}

After the individual preprocessing stages were validated, the complete workflow was executed as an integrated processing pipeline.

This validation focused on verifying the information flow between the different components.

In particular, the output of the automotive event extraction stage was provided to the filtering stage, and the resulting filtered context was subsequently used by the context builder.

This test confirmed that the preprocessing stages could operate together as a complete workflow.

\subsection{Result Validation}

The generated results were inspected to verify their consistency and usefulness for downstream processing.

The validation focused on ensuring that:

\begin{itemize}
    \item relevant communication information was retained;
    \item unnecessary or repetitive information was reduced;
    \item intermediate outputs could be transferred between processing stages;
    \item the final output provided a structured communication context.
\end{itemize}

The validation process demonstrated the importance of evaluating not only the execution of individual components but also the consistency of the information passed between the different stages of the workflow.

\section{Challenges Encountered and Solutions}

Several challenges were encountered during the development of the preprocessing workflow.

One of the main challenges was defining a clear separation between the responsibilities of the different components. Early design decisions could lead to processing logic being distributed between the agent workflow and the tool definitions.

This issue was addressed by simplifying the implementation structure. The specialized tools were maintained as the interface for the agent workflow, while the main processing logic was organized within the core processing layer.

Another challenge concerned the filtering stage. The project objective required unnecessary and redundant traffic to be reduced, but the filtering process also needed to preserve information required for downstream analysis.

This was addressed by treating filtering as a context-preserving operation rather than a simple packet removal mechanism. The objective of the stage is therefore to reduce unnecessary information without blindly discarding communication data that may contribute to the interpretation of system behavior.

A further challenge was ensuring that the output of each tool could be used effectively by the following stage. The workflow was therefore organized around a progressive information transformation process in which each preprocessing operation receives information generated by the previous stage.

Finally, the use of multiple software components required the architecture to remain understandable and maintainable. The resulting implementation was simplified by clearly separating the core processing logic, tool definitions, and agent workflow.

\section{Chapter Summary}

This chapter presented the implementation and validation of the proposed PCAP preprocessing system.

The implementation follows a modular organization that separates the agent workflow, specialized preprocessing tools, and underlying processing logic. The system performs three main preprocessing operations: automotive event extraction, noise filtering and context compression, and context building and report generation.

The workflow uses a shared state to maintain information during execution and coordinates the different processing stages through dedicated nodes.

The implementation was validated through individual tool testing and integrated end-to-end workflow testing. These validation activities focused on verifying the correct transformation and transfer of information between the different preprocessing stages.

The implemented system establishes a preprocessing pipeline capable of progressively transforming communication information into a reduced and structured context suitable for downstream AI-based analysis.